% !TeX root = ../教材题库_全部条目_转写工作稿.tex
% Source: XJ-B1-C06-RV-P264-Y0406
% 第6章 / 章末复习 复习题六 / 教材 p.264 / PDF 第271页
% OCR未校对：公式与文字均可能有误；未完成数学TeX转写；禁止直接发给学生
% 原题及相邻原解截图（核对用）：../题目原图/XJ-B1-C06-RV-P264-Y0406-01.webp
% 本文件只在转写完成后作为题干使用；OCR 条目尚不是可靠题干。
\begin{bankitem}{复习题六 第1题}{第6章 / 章末复习 复习题六 / 教材 p.264 / PDF 第271页}{XJ-B1-C06-RV-P264-Y0406}{OCR未校对：公式与文字均可能有误；未完成数学TeX转写；禁止直接发给学生}
\begin{ocrdraft}
[教材p.264 / PDF第271页]
\par
1.为了解北京市初中学业水平考试中数学成绩的情况，从参加考试的学生中
\par
随机抽查了1000名学生的数学成绩进行统计分析．在这个问题中，下列说法正
\par
确的是（）
\par
（A)总体指的是该市参加学业水平考试的全体学生
\par
(B)个体指的是1佣0名学生中的每一名学生
\par
(C)样本容量指的是1000名学生
\par
（D)样本是指1000名学生的数学学业水平考试成绩
\end{ocrdraft}
\end{bankitem}
